\section{All original Rees coordinate sections on the actual fraction chart}
\begin{definition}[The actual original fraction map's coordinate-chart factor]
\label{def:whole-original-fraction-projective-factor}
\uses{lem:whole-original-fraction-ratio-pivot,lem:whole-original-fraction-ratio-span,def:fc-generated-tuple-twist-frame}
\lean{EnvelopeProperty.BlowupHomogeneousTautologicalLine.originalFractionProjectiveChartFactor}
\leanok
Let $A$ be a domain, $r:R\to A$, $I\subset A$, and $g_0,\ldots,g_m$ the original
specified generators of $I$. Fix an original $i$ with $g_i\ne0$, put
$S=A[I/g_i]$, $a_j=g_j/g_i\in S$, and $r'=(A\to S)r$.
Let $q:\operatorname{Bl}_I(\operatorname{Spec}A)\to\mathbf P_R^m$ be the original
Rees morphism and $c_i:\operatorname{Spec}S\to\operatorname{Bl}_I(\operatorname{Spec}A)$
its actual native fraction inclusion. Define $t_i:\operatorname{Spec}S\to D_+(X_i)$
by the original ratio-tuple chart factor with coefficients $r'$ and entries $a_j$.
\end{definition}
\begin{proof}
The proved original ratio-tuple statements give $a_i=1$ and show that its
entries span the unit ideal. Thus its native coordinate-chart factor is defined. This construction uses the same original $i$th coordinate
open in the same original target $\mathbf P_R^m$.
\end{proof}
\begin{proposition}[The original chart triangle]
\label{prop:whole-original-fraction-projective-factor-triangle}
\uses{def:whole-original-fraction-projective-factor,prop:whole-original-rees-fraction-ratio-map}
\lean{EnvelopeProperty.BlowupHomogeneousTautologicalLine.originalFractionProjectiveChartFactor_comp}
\leanok
Writing $b_i:D_+(X_i)\to\mathbf P_R^m$ for the original coordinate-open inclusion,
one has $t_i b_i=c_iq$ as entire scheme morphisms.
\end{proposition}
\begin{proof}
The native ratio-tuple chart factor followed by $b_i$ is its entire tuple
morphism. The previously proved original Rees fraction equality identifies
that tuple morphism with $c_iq$. Transitivity gives the claimed actual original
triangle; it is not an input certificate.
\end{proof}
\begin{proposition}[The original projective pivot section is a basis]
\label{prop:whole-original-fraction-projective-pivot-basis}
\uses{prop:whole-original-fraction-projective-factor-triangle,thm:fc-coordinate-section-own-chart-frame,thm:fc-line-section-factorization}
\lean{EnvelopeProperty.BlowupHomogeneousTautologicalLine.originalFractionCoordinateSection_isIso}
\leanok
Fix any original native twist pivot $k$. Let $\sigma_j$ be the original global
section of $\mathcal O(1)$ defined by the original coordinate $X_j$, using the
same twist pivot $k$. Then $(c_iq)^*\sigma_i$ is an isomorphism from
$\mathcal O_{\operatorname{Spec}S}$ to its entire actual pullback line.
\end{proposition}
\begin{proof}
On the entire original coordinate open $D_+(X_i)$, the same original section
$b_i^*\sigma_i$ is an isomorphism, as proved by its actual scalar equal to one.
Pulling back through the actual $t_i$ preserves this isomorphism. The native
section composition comparison and the proved triangle $t_i b_i=c_iq$
identify this with $(c_iq)^*\sigma_i$, proving the assertion on its actual whole
target. No basis or section-matching certificate is supplied.
\end{proof}
\begin{proposition}[All original projective sections retain the original ratios]
\label{prop:whole-original-fraction-projective-coordinate-ratios}
\uses{prop:whole-original-fraction-projective-factor-triangle,thm:fc-coordinate-section-chart-ratio,thm:fc-line-section-factor-scalar-relation,thm:fc-generated-tuple-twist-section-scalars}
\lean{EnvelopeProperty.BlowupHomogeneousTautologicalLine.originalFractionCoordinateSection_ratio}
\leanok
For every original index $j$ and every native twist pivot $k$,
\[
(c_iq)^*\sigma_j=\bigl(g_j/g_i\bigr)\,(c_iq)^*\sigma_i
\]
as whole module morphisms on the actual native fraction spectrum. Here the
scalar is the actual element $a_j\in S$, transported to
$\Gamma(\operatorname{Spec}S,\mathcal O)$ by the native spectrum section isomorphism.
\end{proposition}
\begin{proof}
On the original coordinate open $D_+(X_i)$ the whole original coordinate-section
relation is $b_i^*\sigma_j=(X_j/X_i)b_i^*\sigma_i$, with its actual native
coefficient. Pull this relation back along $t_i$, using the native scalar
pullback theorem and the proved actual triangle. The native tuple-factor
coefficient formula identifies $t_i^*(X_j/X_i)$ with precisely the original
$a_j=g_j/g_i\in S$, under the original spectrum-section isomorphism. Substitution
gives the claimed whole original relation. This proves all original coordinate
identities from the original data; it assumes no section matching.
\end{proof}
